\begin{itemize}
    \item 단순한 방법에서 시작할 수 있을까? (Brute Force)
    \item 내가 문제를 푸는 과정을 수식화할 수 있을까? (예제를 직접 해결해보면서)
    \item 문제를 단순화할 수 없을까?
    \item 그림으로 그려볼 수 있을까?
    \item 수식으로 표현할 수 있을까?
    \item 문제를 분해할 수 있을까?
    \item 뒤에서부터 생각해서 풀 수 있을까?
    \item 순서를 강제할 수 있을까?
    \item 특정 형태의 답만을 고려할 수 있을까? (정규화)
    \item \textbf{구간을 통째로 가져간다:} 플로우 + 적당한 자료구조 $(i, i + 1, k, 0)$, $(s, e, 1, w)$, $(N, T, k, 0)$
    \item \textbf{$a = b$:} $a$만 움직이기, $b$만 움직이기, 두 개 동시에 움직이기, 반대로 움직이기
    \item 말도 안 되는 것들을 한 번은 생각해보기 / ``당연하다고 생각한 것'' 다시 생각해보기
    \item \textbf{확률:} DP, 이분 탐색 (NYPC 2019 Finals C)
    \item \textbf{최대/최소:} 이분 탐색, 그리디 (Prefix 고정, Exchange Argument), DP (순서 고정)
    \item 한 가지 풀이에 지나치게 시간을 많이 쓰지 말자. 다른 문제도 한번 읽어보고, 차분히 숨을 돌리며 다른 풀이도 생각해보자.
    \item 접근조차 되지 않는 문제는 어떤 제한을 제거했을 때 문제가 쉬워지는지 판단해볼 필요가 있다.
    \item 억지로 문제를 뚫는 것보다 적절한 관찰을 통하여 좋은 풀이를 작성해내는 것이 더 효과적일 수 있다.
    \item Sometimes, a seemingly naive $O(N^2)$ might turn out to be $O(N \log N)$, or even $O(N)$.
\end{itemize}